%% ============================================================
%% CONCLUSIONS -- v2, restructured into four paragraphs.
%% Replaces the single block in the skeleton. NO SUBHEADINGS (editor requirement).
%% Pattern follows the Applied Energy exemplar: what was developed and tested; what was
%% found; what it means for model selection; scope and extension.
%%
%% The third paragraph now LEADS with the loss-number criterion. It is the most portable
%% result in the paper -- a decision a planner can make before building anything -- and in
%% the previous version it was buried mid-sentence in the second half of a long block.
%% ============================================================

This paper set out to determine which modelling choices actually change a district-heating
dispatch decision, rather than which change a reported objective value. To do so it
separates two quantities that the fidelity literature has treated as one (estimation
bias, the difference in optimised cost between formulations, and decision regret, the cost
of executing a formulation's schedule under a common forward model) and evaluates both on
a ladder in which each rung adds exactly one phenomenon. A copperplate supplied with the
aggregate loss it cannot itself compute serves as a control condition, so that loss
representation and spatial resolution can be switched independently and their contributions
separated exactly rather than estimated. The design was applied to one industrial network
with measured data and to a complete 135-cell balanced factorial of synthetic networks.

What a simplified model must capture is loss visibility, not spatial topology. The
copperplate-to-baseline cost gap decomposes to machine precision into a \textbf{95.8}\,\%
loss main effect and a \textbf{4.7}\,\% spatial-topology main effect, and the same holds
across the synthetic factorial. The distinction is decision-relevant, not book-keeping: a
copperplate understates cost by \textbf{15.1}\,\% yet, executed, costs \textbf{46.1}\,\%
more to operate---opposite signs. The schedules show why: blind to losses, it provisions the
heat pump for too small a demand, running it some 246 fewer hours at ten points less share,
and tops up the shortfall at the margin in winter hours it never planned for. A lumped
per-network loss correction captures most of the bias, but only where its coefficient is
known: frozen and carried to another network it mis-estimates the loss burden by tens of
percentage points, and even where it reproduces the cost it runs a materially different
generation mix. At the opposite extreme---the demand side resolved to every transmission
station and service lateral, on manufacturer data with an independent nonlinear
cross-check---hydraulic detail changes cost by under one percent and decisions by essentially
nothing.

For model selection this yields a criterion that can be applied before any model is built.
The share of cost a copperplate misses is $b = \lambda/(1+\lambda)$, where the loss number
$\lambda$ is the ratio of annual network loss to annual demand and follows from the pipe
inventory alone. Across all 136 networks examined this parameter-free relation reproduces
the measured burden with $R^2 = \textbf{0.87}$. The screening decision (copperplate,
calibrated adder, or resolved nodes) is read off from $\lambda$ alone, before any model is
built. Where a model must be resolved, effort belongs in the loss representation before
the network graph and well before the hydraulics; and fidelity should be claimed only to the
resolution at which it can be verified, which for operators metering at billing grade is the
node or zone level rather than the substation.

These conclusions are bounded by the regime that produced them, and the boundaries are
specific rather than decorative. They hold for radial, centrally-supplied networks with
fixed capacities, a prescribed heating curve and hourly deterministic dispatch. Each of
those assumptions closes a channel through which extended physics could act, so the nulls
reported here are conservative lower bounds: relaxing any of them can give the physics more
room to matter, never less, as the supply-temperature study demonstrates by opening one of
them and turning hydraulics from a negligible cost into the binding constraint. Ring and
bidirectional topologies require a signed-flow formulation absent here; distributed
generation remains an open question because every case studied is centrally supplied; and
post-upgrade measurement of the new assets' dispatch is the natural next validation step.
The sharpest boundary is to the companion study: whether network resolution changes not only
how a network is dispatched but where new capacity is sited, under joint
dispatch-and-sizing, is a question this paper deliberately does not settle.
